% Revision is a diagnostic loop, not an additional feedback category.
\begin{tikzpicture}[
  x=1mm,y=1mm,font=\figurefont\fontsize{8}{10}\selectfont,
  box/.style={draw=BookRule,fill=BookPaper,text width=32mm,
    minimum height=12mm,align=center,inner sep=2mm,line width=.7pt},
  flow/.style={-{Stealth[length=2mm]},draw=BookMuted,line width=.8pt},
  lab/.style={fill=white,inner sep=1mm,font=\figurefont\fontsize{7}{9}\selectfont}
]
  \node[box,fill=FigureInput!10] (goal) at (20,36) {目标要求与训练代理};
  \node[box,fill=FigureOutput!10] (check) at (81,36) {独立检测};
  \node[box] (diagnose) at (143,36) {根因判断};
  \node[box,fill=FigureModel!10] (revision) at (143,0)
    {修订要求／模型／规则\\或限制策略更新};
  \node[box] (behavior) at (81,0) {更新后新行为};
  \node[box] (version) at (20,0) {版本对应与反例记录};
  \draw[flow] (goal) -- node[lab,above] {要求与检查对象} (check);
  \draw[flow] (check) -- node[lab,above] {具体反例} (diagnose);
  \draw[flow] (diagnose) -- node[lab,right] {待修补缺口} (revision);
  \draw[flow] (revision) -- node[lab,above] {新策略} (behavior);
  \draw[flow] (behavior) -- node[lab,left] {重新检测} (check);
  \draw[flow] (revision.south) -- ++(0,-12) -| node[lab,near start,below]
    {保留改变了什么} (version.south);
  \draw[flow] (version) -- node[lab,left] {修订后的约定} (goal);
\end{tikzpicture}
